% Avsnitt 18.9, ur lectures/L18/README.md:s SPI-genomgång och den utdelade bridge/spi_slave.vhd.

\section{SPI från vågformen och uppåt}
\label{c18:sec:spi}

Registerbanken presenterar register. Något måste nå dem utifrån, och det är SPI. Det här avsnittet
tar transporten från vågformen och uppåt, och slutar i den utdelade \code{spi_slave}, som du läser
snarare än skriver.

\subsection*{Läge 0, MSB först, SS som ramning}
Projektet kör \textbf{SPI läge 0} (CPOL~=~0, CPHA~=~0) rakt igenom:
\begin{itemize}
  \item \code{SCK} vilar lågt. Mastern genererar den; slaven producerar aldrig en klocka.
  \item Båda sidor \textbf{samplar på den stigande flanken}. \code{MISO} ändras därför på den
    \textbf{fallande} flanken, så att värdet hunnit stabilisera sig till nästa stigande.
  \item \textbf{MSB först}, i varje byte.
  \item \code{SS} är aktiv låg och är \textbf{ramningssignalen}: en transaktion per låg period.
    Mastern måste höja \code{SS} mellan transaktioner.
  \item \code{SCK} går i högst 1~MHz. Skälet står i \secref{c18:sub:oversample}.
\end{itemize}

Rita vågformen för hand innan du läser vidare: \code{SS} som faller, åtta stigande \code{SCK}-flanker
med byten \code{0xA5} på \code{MOSI}, mest signifikant först, och \code{SS} som stiger. Det är hela
bilden, och Övning~\ref{c18:ex:waveform} bygger en femtransaktionsbyte av den.

\subsection{Varför slaven inte klockar på SCK}
\label{c18:sub:oversample}

Den frestande läsningen är att \code{SCK} är en klocka, så gör den till \emph{din} klocka: klocka
skiftregistret på den och var klar. Låt bli, av två skilda skäl.

\textbf{Det första är metastabilitet}, och det är \chapref{c4:ch}:s regel oförändrad. \code{SCK},
\code{MOSI} och \code{SS} kommer alla in på pinnar som drivs av en mikrokontroller med sin egen
kristall. De är asynkrona mot \code{clock} i precis den mening \chapref{c4:ch} avser, och varje en av
dem måste genom två vippor innan någon logik får titta på den.

\textbf{Det andra är vad en design med två orelaterade klockor kostar i en FPGA.} En andra klockdomän
betyder en andra uppsättning tidsvillkor, ett verkligt klockdomänkors på varje signal som passerar
mellan dem, och ett syntesverktyg som inte längre kan analysera designen som en helhet. För ett
gränssnitt som går i 1~MHz mot en väv som går i 50 är det priset onödigt.

Svaret är \textbf{översampling}: synkronisera in \code{SCK} i 50 MHz-domänen och \emph{detektera dess
flanker} i stället för att klocka på den. Vid 1~MHz är det minst 50 systemcykler per SPI-bit, alltså
en storleksordnings marginal. Protokollet garanterar därför korrekt funktion upp till 1~MHz; snabbare
kan fungera, och lovas medvetet inte.

Det kostar några 50 MHz-cykler per flank, och det är därifrån protokollets tidskrav kring \code{SS}
kommer: minst 60~ns (tre cykler) före den första \code{SCK}-flanken, mellan transaktioner, och efter
den sista flanken. Bilaga~\ref{app:protocol} har hela tabellen.

\subsection{Att läsa \code{spi_slave.vhd}}
\label{c18:sub:slave}

\code{spi_slave} delas ut i stället för att skrivas: den är transport, inte bokens ämne. Men du måste
kunna läsa den, och \chapref{c20:ch} förväntar sig att du kan det.

\vhdlfile{bridge/spi_slave.vhd}

Fyra saker är värda att läsa noga:

\paragraph{Posten \code{sync_t} har tre fält, men bara två är synkroniseringssteg}
\code{s1} och \code{s2} är dubbelvippsynkroniseraren; \code{s3} är förra cykelns \code{s2}, och finns
bara för att en flank ska gå att detektera genom att jämföra de två. \code{sclk} och \code{ss}
används för sina \emph{flanker} och behöver alla tre; \code{mosi} läses bara vid en flank som redan
detekterats på \code{sclk}, så den behöver de två synkroniseringsstegen och ingen historik.

\paragraph{De två kedjorna resettar till olika värden}
\code{sclk_s} nollställs till idel nollor och \code{ss_s} till idel \textbf{ettor}. Varje kedja
resettar till den nivå dess ledning vilar på, så att modulen kommer ur reset och tror på det som
faktiskt är sant om en ledig buss: \code{sclk} vilar lågt i läge 0, och \code{ss} är aktiv låg och
vilar \textbf{högt}. Att nollställa \code{ss}-kedjan till \code{'0'} skulle få \code{ss_active} att
läsa \code{'1'} ur reset och bryggan att se en ram som ingen startat. Det är exakt samma regel som
\code{button_sync} i Övning~\ref{c4:ex:buttonsync} finns till för att upprätthålla.

\paragraph{En avvald slav ignorerar SPI-ledningarna helt}
Hanteringen av \code{SCK}-flanker ligger inuti \code{else}-armen i ett test på \code{ss_s.s2}.
\code{SCK} tillhör den slav mastern adresserar; för den här slaven är den en asynkron ingång som
kommer in över en ledning på ett kopplingsdäck, och en störning på den ledningen ser likadan ut som
en flank. Utan den grindningen skulle grenarna nedanför skifta in \code{MOSI} och pulsa
\code{rx_valid} på flanker som anländer medan \code{ss} är hög, och \code{spi_reg_bridge} skulle
tappa bytejusteringen för varje transaktion därefter.

\paragraph{\code{ss_falling} är ett \code{elsif} före SCK-grenarna}
Ordningen spelar roll vid exakt en klockflank: den där \code{SS} nyss föll \emph{och} en stigande
\code{SCK}-flank detekteras samtidigt. Genom att \code{ss_falling}-armen kommer först kan den
\code{SCK}-flanken inte stega biträknaren förbi ramens början och kasta om varje byte i den.

\code{spi_slave} lämnar över en färdig byte i taget på \code{rx_data}, med \code{rx_valid} som en
puls, och tar nästa byte att sända på \code{tx_data}. Den vet ingenting om register.
\Chapref{c19:ch} bygger \code{spi_reg_bridge}, lagret som gör de byteströmmarna till
femtransaktionsbytes och till registeråtkomster mot banken du just skrivit.
